(1) 在旋转框架与单激发子空间 \(|e,0\rangle,|g,1\rangle\) 中，振幅可写 \(\dot c_e=-igc_a\)、\(\dot c_a=-(\kappa/2-i\Delta)c_a-igc_e\)。这里 \(\kappa\) 是腔光子数能量衰减率，故振幅的损失项为 \(-\kappa c_a/2\)；背景发射体辐射暂时忽略。

(2) 当 \(g\ll\kappa\) 且腔模快速弛豫，近似 \(c_a\simeq-igc_e/(\kappa/2-i\Delta)\)。代回 \(\dot c_e\) 并取实部得激发态布居衰减率 \(\boxed{\Gamma_{\rm cav}=g^2\kappa/[(\kappa/2)^2+\Delta^2]}\)。在 \(\Delta=0\) 时 \(\boxed{\Gamma_{\rm cav}=4g^2/\kappa}\)，大失谐的增强以Lorentz因子 \([1+4\Delta^2/\kappa^2]^{-1}\) 下降。

(3) 有效模体积 \(V_{\rm eff}\) 处的单光子电场强度满足 \(E_{1\rm ph}=\sqrt{\hbar\omega_c/(2\varepsilon_0n^2V_{\rm eff})}\)，耦合 \(g=|\mathbf d\cdot\mathbf E_{1\rm ph}|/\hbar\)。均匀背景发射率 \(\gamma_0=n\omega_e^3|\mathbf d|^2/(3\pi\varepsilon_0\hbar c^3)\)。共振且位置/方向理想匹配时 \(\boxed{F_P=\Gamma_{\rm cav}/\gamma_0=\frac{3}{4\pi^2}(\lambda_0/n)^3\frac Q{V_{\rm eff}}}\)，其中 \(\lambda_0=2\pi c/\omega_c\)。

(4) 若其他背景通道仍以 \(\gamma_0\) 辐射，进入腔模的分支比为 \(\boxed{\beta=\Gamma_{\rm cav}/(\Gamma_{\rm cav}+\gamma_0)}\)，共振时等于 \(F_P/(1+F_P)\)。若 \(g\) 与 \(\kappa\) 同量级，则腔—发射体形成可分辨的混合Rabi态，不能再简单绝热消去腔光子并声称有单一指数寿命。
